\scanpage{61}
\margin{$-,6$}
$1^\circ$) Non \unsure{tordu|tordus} tous deux \quad $e_1\unsure{-|\to}e_1\quad e_2-e_3,\ e_3-e_2$

symétrie \unsure{autour|suivant} plan diagonal (\uline{ordre 2})

\margin{$+,6$}
$2^\circ$) 1 tordu, 2 non tordu \quad $e_1-{-e_1},\ e_2-e_3,\ e_3-e_2$

symétrie autour axe-arête \inserted{\unsure{conjugué|composée} au produit du préc. par antip.} (\uline{ordre 2})

\margin{$+,6$}
$3^\circ$) 1 non tordu, 2 tordu \quad $e_1-e_1,\ e_2\to e_3,\ e_3-{-e_2}$ (ordre 4)

rotation d'un quart de tour autour de l'axe perpendiculaire à \unsure{une|deux} face (axe-face)

\margin{$-,6$}
$4^\circ$) 1 tordu, 2 tordu \quad $e_1-{-e_1},\ e_2-e_3,\ e_3-{-e_2}$ (ordre 4)

(conjugué au précédent, suivi de \deleted{\unsure{symétrie}} \inserted{antipodisme})

c) Autom décomposables \unsure{à|.} trois cycles (déterminé par $\alpha_1^+$ et $\alpha_1^-$, $\alpha_1^++\alpha_1^-=3$)

\margin{$+,1$}
$1^\circ$) $\alpha_1^+=3$\qquad identité

\margin{$-,3$}
$2^\circ$) $\alpha_1^+=2$\qquad $e_1\to e_1,\ e_2-e_2,\ e_3\to -e_3$\qquad symétrie par rapport à plan \deleted{\unsure{d'axe}} coord.

\margin{$+,3$}
$3^\circ$) $\alpha_1^+=1$\qquad $e_1\to e_1,\ e_2-{-e_2},\ e_3-{-e_3}$\qquad sym. \unsure{par r.|par rapport} à axe coord

\margin{$-,1$}
$4^\circ$) $\alpha_1^+=0$\qquad antipodisme

En tout il y a $1\deleted{\unsure{4}}0$ classes \unsure{,|?}
\[
(10=2+4+4)
\]
\noindent\rule{\linewidth}{0.4pt}

\uline{Prop} Soit $u$ un autom de $C$ tel que pour tt sommet $x$, on ait $ux=\pm x$. Alors \inserted{(si $n\geq3$)} $u$ est l'identité ou l'antipodisme.

\unsure{Donc|D'où}
\[
\operatorname{Aut}(C)/(\pm1)\underset{\text{\uuline{inj}}}{\lhook\joinrel\longrightarrow}\operatorname{Aut}(\Phi_0(C)/\sigma)
\]
donc si $C$ de dim impaire
\[
\operatorname{Aut}^+(C)\unsure{\lhook\joinrel\relbar|\hookrightarrow}\operatorname{Aut}(\unsure{\Phi_0(C)|\Phi(C)})
\]
% archive sheet number: 60
